\documentclass[svgnames]{beamer}
\mode<presentation>
\usefonttheme{serif}
\usecolortheme{dove}
\useinnertheme{rounded}
%\useoutertheme{smoothbars}
\setbeamercolor{item projected}{fg=black}
\setbeamertemplate{navigation symbols}{}

\usepackage[english]{babel}
\usepackage[latin1]{inputenc}
\usepackage{times}
\usepackage{amsthm,amssymb,amsmath,graphicx}
\usepackage{color}
\usepackage{gastex}
\usepackage{tikz}

\renewcommand{\AA}{\mathcal{A}}

\newcommand{\PSPACE}{\mathrm{PSPACE}}
\newcommand{\BB}{\mathcal{B}}
\newcommand{\CC}{\mathcal{C}}
\newcommand{\JJ}{\mathcal{J}}
\newcommand{\DD}{\mathcal{D}}
\newcommand{\KK}{\mathcal{K}}
\newcommand{\LL}{\mathcal{L}}
\newcommand{\HH}{\mathcal{H}}
\newcommand{\GG}{\mathcal{G}}
\newcommand{\RR}{\mathbb{R}}
\newcommand{\NN}{\mathbb{N}}
\newcommand{\QQ}{\mathbb{Q}}
\newcommand{\MM}{\mathcal{M}}
\newcommand{\MMAA}{\MM_\AA}
\newcommand{\tr}[1]{\langle #1 \rangle}
\newcommand{\prob}[1]{\mathbb{P}_{#1}}
\newcommand{\stab}[1]{S(#1)}
\newcommand{\set}[1]{\{ #1 \}}
\newcommand{\val}[1]{\text{val}(#1)}
\newcommand{\merge}{\textrm{merge}}
\newcommand{\chck}{\textrm{check}}
\newcommand{\apply}{\textrm{apply}}
\newcommand{\finish}{\textrm{finish}}
\newcommand{\wait}{\textrm{wait}}
\newcommand{\delay}{\textrm{Delay}}
   
\title{Deciding the value $1$ problem \\ for probabilistic leaktight automata}
\author{Nathana\"el Fijalkow, \\
joint work with Hugo Gimbert and Youssouf Oualhadj}
\institute{LIAFA, CNRS \& Universit\'e Denis Diderot - Paris 7, France\\
\textbf{nath@liafa.jussieu.fr}}
\date{November 25th, 2011}
\subtitle{S\'eminaire Automates}

\AtBeginSection[]
{
  \begin{frame}<beamer>{Outline}
    \tableofcontents[currentsection,currentsubsection]
  \end{frame}
}
\AtBeginSubsection[]
{
  \begin{frame}<beamer>{Outline}
    \tableofcontents[currentsection,currentsubsection]
  \end{frame}
}

\begin{document}

\begin{frame}
  \titlepage
\end{frame}

\begin{frame}<beamer>{Outline}
  \tableofcontents
\end{frame}

\section{The value $1$ problem for probabilistic automata}

\subsection{Definitions}

\begin{frame}{Probabilistic automata (Rabin, 1963)}
\begin{figure}
\begin{center}
\begin{picture}(20,40)(0,-10)
	\gasset{Nw=6,Nh=6}

  	\node[Nmarks=i,iangle=0](L1)(15,15){$1$}
  	\node(L2)(0,30){$2$}
  	\node[Nmarks=r](L3)(0,0){$3$}

  	\drawedge(L1,L2){$b$}
  	\drawedge[curvedepth=-5,ELside=r](L1,L3){$a,.4$}
  	\drawedge[curvedepth=-5,ELside=r](L3,L1){$b$}
	\drawloop(L1){$a,.6$}
	\drawloop[loopangle=135](L2){$a,b$}
	\drawloop[loopangle=215](L3){$a$}
\end{picture}
\end{center}
\end{figure}
$$\prob{\AA} : A^* \rightarrow [0,1]$$
\end{frame}

\begin{frame}{Cutpoint and value}
Fix $0 < \lambda \leq 1$, define:
$$L_\lambda = \set{w \mid \prob{\AA}(w) \geq \lambda}.$$

\pause
$\lambda$ is isolated if there exists $\delta > 0$ such that for all $w \in A^*$, we have
$$|\prob{\AA}(w) - \lambda| \geq \delta$$

\pause
\begin{theorem}[Rabin, 1963]
If $\lambda$ is isolated, then $L_\lambda$ is a regular language.
\end{theorem}

\end{frame}

\subsection{Deciding the isolation problem}

\begin{frame}{The isolation problem}

Fix $0 \leq \lambda \leq 1$, the isolation problem is:
\begin{itemize}
	\item[] \textbf{Instance:} a probabilistic automaton $\AA$
	\item[] \textbf{Question:} is $\lambda$ isolated in $\AA$?
\end{itemize}

\pause
For $0 < \lambda < 1$, Bertoni showed that this is undecidable (in 1974)!
\end{frame}

\begin{frame}{The value $1$ problem}
For $\lambda = 1$ the isolation problem can be formulated as:
``are there words accepted by $\AA$ with probability arbitrarily close to $1$''.
\pause

Equivalently, define $\val{\AA} = \sup_w \prob{\AA}(w)$,
then the problem is:
\begin{center}
``$\val{\AA} \stackrel{?}{=} 1$''.
\end{center}

\pause
\begin{theorem}[Gimbert, Oualhadj, 2010]
The value $1$ problem is undecidable.
\end{theorem}
\end{frame}

\begin{frame}{An intuition}

\begin{figure}
\begin{center}
\begin{picture}(60,35)(0,0)
	\gasset{Nw=7,Nh=7}

  	\node[Nmarks=i,iangle=-90](0)(30,15){$0$}
  	\node(L1)(0,10){$L_1$}
  	\node[Nmarks=r](L2)(0,30){$L_2$}
  	\node(R1)(60,10){$R_1$}
  	\node(R2)(60,30){$R_2$}

	\drawloop(0){$a$}

  	\drawedge[curvedepth=5,ELside=l](0,L1){$b,\frac{1}{2}$}
  	\drawedge[curvedepth=5,ELside=l](L1,0){$a,1-x$}
  	\drawedge(L1,L2){$b$}
	\drawloop[loopangle=-135](L1){$a,x$}
	\drawloop[loopangle=90](L2){$a,b$}

  	\drawedge[curvedepth=-5,ELside=r](0,R1){$b,\frac{1}{2}$}
  	\drawedge[curvedepth=-5,ELside=r](R1,0){$a,x$}
  	\drawedge[ELside=r](R1,R2){$b$}
	\drawloop[loopangle=-45](R1){$a,1-x$}
	\drawloop(R2){$a,b$}
\end{picture}
\end{center}
\end{figure}
\begin{center}
has value $1$ if and only if $x > \frac{1}{2}$.
\end{center}
\end{frame}

\begin{frame}{A very restricted case}
\begin{theorem}[Fijalkow, Gimbert, Oualhadj, 2011]
The isolation problem is (still) undecidable if we randomise only on \textbf{one} transition.
\end{theorem}
\end{frame}

\begin{frame}{Sketch of proof (1)}
\begin{figure}
\begin{center}
\begin{picture}(120,30)(7,0)
	\gasset{Nw=5,Nh=5,Nframe=n}

	\put(10,30){$a$}
	\drawline[dash={1.5}0,AHnb=0](10,-2)(10,27)
	\put(30,30){$b$}
	\drawline[dash={1.5}0,AHnb=0](30,-2)(30,27)
	\put(50,30){$b$}
	\drawline[dash={1.5}0,AHnb=0](50,-2)(50,27)
	\put(70,30){$a$}
	\drawline[dash={1.5}0,AHnb=0](70,-2)(70,27)
	\put(90,30){$b$}
	\drawline[dash={1.5}0,AHnb=0](90,-2)(90,27)
	\put(110,30){$b$}
	\drawline[dash={1.5}0,AHnb=0](110,-2)(110,27)

  	\node(1)(0,15){$q_0$}
  	\node(2)(20,15){$q_1$}
  	\node(3)(40,15){$q_2$}
  	\node(4)(60,25){$r$}
  	\node(5)(60,5){$s$}
  	\node(6)(80,25){$\ldots$}
  	\node(7)(80,12){$\ldots$}
  	\node(8)(80,0){$\ldots$}
  	\node(9)(100,25){$q$}
  	\node(9-10)(100,21){}
  	\node(10)(100,17){$q$}
  	\node(11)(100,5){$\ldots$}
  	\node(12)(100,0){$\ldots$}
  	\node(13)(120,0){$\ldots$}
  	\node(14)(120,3){$\ldots$}
  	\node(15)(120,8){$\ldots$}
  	\node(16)(120,21){$\ldots$}

  	\drawedge(1,2){}
  	\drawedge(2,3){}
  	\drawedge(3,4){}
  	\drawedge(3,5){}
  	\drawedge(4,6){}
  	\drawedge(5,7){}
  	\drawedge(5,8){}
  	\drawedge(6,9){}
  	\drawedge(7,10){}
  	\drawedge(7,11){}
  	\drawedge(8,12){}
  	\drawedge(9-10,16){}
  	\drawedge(11,15){}
  	\drawedge(11,14){}
  	\drawedge(12,13){}

	\put(99,20){$\parallel$}
\end{picture}
\end{center}
\end{figure}
\end{frame}

\begin{frame}{Sketch of proof (2)}
Given $\AA$ reading words from $A^*$, we construct $\BB$ over a new alphabet $B$, with one probabilistic transition,
and a morphism $\widehat{\_} : A^* \rightarrow B^*$ such that:
$$\forall w \in A^*, \prob{\AA}(w) = \prob{\BB}(\widehat{w}).$$

\pause
$$\widehat{a} =
\chck(a,q_0) \cdot * \cdot \apply(a,q_0) \ldots \chck(a,q_{n-1}) \cdot * \cdot \apply(a,q_{n-1}) \cdot \merge.$$

\pause
\begin{figure}
\begin{center}
\begin{picture}(80,20)(0,0)
	\gasset{Nw=5,Nh=5,Nframe=n}

  	\node(q)(0,10){$q$}
  	\node(s)(25,10){$s_*$}
  	\node(s0)(40,20){$s_0$}
  	\node(s1)(40,0){$s_1$}
  	\node(hr0)(70,20){$\widetilde{r_0}$}
  	\node(hr1)(70,0){$\widetilde{r_1}$}
  	\node(r0)(85,20){$r_0$}
  	\node(r1)(85,0){$r_1$}

  	\drawedge(q,s){$\chck(a,q)$}
  	\drawedge(s,s0){$*$}
  	\drawedge[ELside=r](s,s1){$*$}
  	\drawedge[ELside=r](s0,hr0){$\apply(a,q)$}
  	\drawedge(s1,hr1){$\apply(a,q)$}
  	\drawedge[ELside=r](hr0,r0){$\merge$}
  	\drawedge(hr1,r1){$\merge$}
\end{picture}
\end{center}
\end{figure}
\end{frame}

\begin{frame}{Sketch of proof (3)}
%$$\widehat{a} = \begin{array}{l}
%\chck(a,p) \cdot * \cdot \apply(a,p) \cdot \\
%\chck(a,q) \cdot * \cdot \apply(a,q) \cdot \\
%\chck(a,r) \cdot * \cdot \apply(a,r) \cdot \\
%\cdot \merge
%\end{array}$$
\begin{figure}
\begin{center}
\begin{picture}(100,20)(0,0)
	\gasset{Nw=5,Nh=5,Nframe=n}

	\put(10,35){$\AA$}
	\put(10,25){$a$}
	\drawline[dash={1.5}0,AHnb=0](10,-4)(10,22)

	\put(50,35){$\BB$}
  	\node(p)(0,20){$p$}
  	\node(q)(0,10){$q$}
  	\node(r)(0,0){$r$}
  	\node(wp)(20,20){$p$}
  	\node(wr)(20,10){$r$}
  	\node(wq)(20,0){$q$}

  	\drawedge(p,wp){}
  	\drawedge(q,wp){}
  	\drawedge(q,wr){}
  	\drawedge(r,wq){}

	\only<2>{\put(40,25){$\chck(a,p) \cdot * \cdot \apply(a,p)$}}
	\only<3>{\put(40,25){$\chck(a,q) \cdot * \cdot \apply(a,q)$}}
	\only<4>{\put(40,25){$\chck(a,r) \cdot * \cdot \apply(a,r)$}}
	\only<5>{\put(85,25){$\merge$}}
	\drawline[dash={1.5}0,AHnb=0](10,-4)(10,22)

  	\node(pp)(40,20){$p$}
  	\node(qp)(40,10){$q$}
  	\node(rp)(40,0){$r$}

  	\only<2,3,4,5>{
  		\node(wp)(80,20){$\widetilde{p}$}
	  	\drawedge(pp,wp){}
  		}
  	\only<3,4,5>{
  		\node(wr)(80,10){$\widetilde{r}$}
  		\drawedge(qp,wp){}
  		\drawedge(qp,wr){}
  		}
  	\only<4,5>{
  		\node(wq)(80,0){$\widetilde{q}$}
	  	\drawedge(rp,wq){}  	
  		}
  	\only<5>{
  		\node(p2)(100,20){$p$}
  		\node(r2)(100,10){$r$}
  		\node(q2)(100,0){$q$}
	  	\drawedge(wp,p2){}  	
	  	\drawedge(wr,r2){}  	
	  	\drawedge(wq,q2){}  	
  		}
\end{picture}
\end{center}
\end{figure}
\end{frame}

\begin{frame}{Sketch of proof (4)}
$\BB$ is unable to check that a letter $\chck(a,q)$ is actually followed by the corresponding $\apply(a,q)$: 
inbetween, it will go through $s_*$ and ``forget'' the state it was in.
\pause

\begin{figure}
\begin{center}
\begin{picture}(80,28)(0,0)
	\gasset{Nadjust=wh,Nframe=n}

  	\node(q)(0,20){$q$}
  	\node(s)(40,20){$s_*$}
  	\node(s0)(55,30){$s_0$}
  	\node(s1)(55,10){$s_1$}
  	\node(r0)(80,30){$\widetilde{r_0}$}
  	\node(r1)(80,10){$\widetilde{r_1}$}

  	\node(w)(40,5){$w$}
  	\node(q_0)(10,5){to $q_0$}

  	\drawedge(q,s){$\chck(a,q)$}
  	\drawedge[ELside=r](s,w){$*$}
	\drawloop[loopangle=-90](w){}
  	\drawedge(w,q_0){$\finish$}
  	\drawedge[ELside=l](s,s0){$*$}
  	\drawedge[ELside=r](s,s1){$*$}
  	\drawedge(s0,r0){$\apply(a,q)$}
  	\drawedge[ELside=r](s1,r1){$\apply(a,q)$}
\end{picture}
\end{center}
\end{figure}

$$\sup_n \prob{\BB}((\widehat{w} \cdot \finish)^n) = \prob{\AA}(w)$$
\end{frame}

\begin{frame}{Sketch of proof (5)}
Assume $p \in F$, $q \notin F$ 
and $i$ is the initial state of a (deterministic) automaton recognizing $(\widehat{A^*} \cdot \finish)^*$.

\begin{figure}
\begin{center}
\begin{picture}(40,28)(0,0)
	\gasset{Nw=5,Nh=5,Nframe=n}

	\put(9,25){$\widehat{w}$}
	\drawline[dash={1.5}0,AHnb=0](10,-4)(10,22)
	\put(26,25){$\finish$}
	\drawline[dash={1.5}0,AHnb=0](30,-4)(30,22)

  	\node(q_0)(0,10){$q_0$}
  	\node(p)(20,20){$p$}
  	\node(q)(20,10){$q$}
  	\node(w)(20,0){$w$}
  	\node(q_02)(40,0){$q_0$}
  	\node(i)(40,20){$i$}

  	\drawedge(q_0,p){}
  	\drawedge(q_0,q){}
  	\drawedge(q_0,w){}
  	\drawedge(w,q_02){}
  	\drawedge(q,q_02){}
  	\drawedge(p,i){}
\end{picture}
\end{center}
\end{figure}
\end{frame}

\begin{frame}{Our objective}
Define a \textbf{large} and \textbf{interesting} subclass of probabilistic automata for which the value $1$ problem is decidable.
\end{frame}

\section{An algebraic solution to the limitedness problem for distance automata}

\subsection{Taking a step back: weighted automata}

\begin{frame}{Probabilistic automata VS distance automata}
Consider a semiring $(\KK,+,\cdot)$.
An automaton computes in the semiring $\KK$ if
$\val{w} = \sum \set{\Pi(\rho) \mid \rho \textrm{ is a run over } w}$.
\pause

\begin{itemize}
	\item Classical automata compute in the boolean semiring.
	\item Probabilistic automata compute in $(\RR,+,\cdot)$ (there is a catch here).
	\item Distance automata compute in the tropical semiring $(\NN \cup \set{\infty},\min,+)$.
	Here is an example:
\end{itemize}

\begin{figure}
\begin{center}
\begin{picture}(40,15)(0,0)
	\gasset{Nw=6,Nh=6}

  	\node[Nmarks=i,iangle=-90](p)(0,0){$p$}
  	\node[Nmarks=ir,iangle=-90](q)(20,0){$q$}
  	\node[Nmarks=r](r)(40,0){$r$}

	\drawloop(p){$a,b:0$}
  	\drawedge(p,q){$b:0$}
	\drawloop(q){$a:1$}
  	\drawedge(q,r){$b:0$}
	\drawloop(r){$a,b:0$}
\end{picture}
\end{center}
\end{figure}
\end{frame}

\begin{frame}{The value $1$ problem VS the limitedness problem}
The value $1$ problem for probabilistic automata is: 
\begin{center}
``are there words accepted with probability arbitrarily close to $1$?''.
\only<2>{\color{red}{undecidable}}
\end{center}

\vskip3em

The \textbf{un}limitedness problem for distance automata is: 
\begin{center}
``are there words with arbitrarily high value?''.

\pause
\color{red}{decidable} 
(Hashiguchi, 1988)
\end{center}
\end{frame}

\subsection{Leung's algorithm}

\begin{frame}{Weighted automata using algebra (Sch\"utzenberger)}
\begin{figure}
\begin{center}
\begin{picture}(40,10)(0,0)
	\gasset{Nw=6,Nh=6}

  	\node[Nmarks=i,iangle=-90](p)(0,0){$p$}
  	\node[Nmarks=ir,iangle=-90](q)(20,0){$q$}
  	\node[Nmarks=r](r)(40,0){$r$}

	\drawloop(p){$a,b:0$}
  	\drawedge(p,q){$b:0$}
	\drawloop(q){$a:1$}
  	\drawedge(q,r){$b:0$}
	\drawloop(r){$a,b:0$}
\end{picture}
\end{center}
\end{figure}

$$\tr{a} = 
\left(\begin{array}{ccc}
0 & \infty & \infty \\
\infty & 1 & \infty \\
\infty & \infty & 0
\end{array}\right)
\qquad
\tr{b} = 
\left(\begin{array}{ccc}
0 & 0 & \infty \\
\infty & \infty & 0 \\
\infty & \infty & 0
\end{array}\right)$$

\only<1>{
$$I = \left(\begin{array}{ccc} 0 & 0 & \infty \end{array}\right)
\qquad
F = \left(\begin{array}{c} \infty \\ 0 \\ 0 \end{array}\right)$$}

\only<2>{
$$I \cdot \tr{aaabaa} \cdot F = 
\left(\begin{array}{ccc} 0 & 0 & \infty \end{array}\right) \cdot
\left(\begin{array}{ccc}
0 & 2 & 2 \\
\infty & \infty & 3 \\
\infty & \infty & 0
\end{array}\right) \cdot
\left(\begin{array}{c} \infty \\ 0 \\ 0 \end{array}\right) = 2$$

$$\left\{
\begin{array}{cc}
k \in \NN & \textrm{best run has value } k\\
\infty & \textrm{no run}
\end{array}\right.$$}
\end{frame}

\begin{frame}{Towards Leung's algorithm: $\sharp$-expressions}
\begin{figure}
\begin{center}
\begin{picture}(70,10)(0,0)
	\gasset{Nw=6,Nh=6}

  	\node[Nmarks=i,iangle=-90](p)(0,5){$p$}
  	\node[Nmarks=ir,iangle=-90](q)(20,5){$q$}
  	\node[Nmarks=r](r)(40,5){$r$}

	\drawloop(p){$a,b:0$}
  	\drawedge(p,q){$b:0$}
	\drawloop(q){$a:1$}
  	\drawedge(q,r){$b:0$}
	\drawloop(r){$a,b:0$}

  	\node[Nmarks=ir,iangle=180](s)(70,5){$s$}
	\drawloop[loopangle=90](s){$a:0$}
	\drawloop[loopangle=-90](s){$b:1$}
\end{picture}
\end{center}
\end{figure}

$$\val{(a^n \cdot b)^n \cdot a^n} = n$$
\pause
An unlimitedness witness is $(a^\sharp \cdot b)^\sharp \cdot a^\sharp$.
\end{frame}

\begin{frame}{Towards Leung's algorithm: stabilization}
\begin{figure}
\begin{center}
\begin{picture}(70,10)(0,0)
	\gasset{Nw=6,Nh=6}

  	\node[Nmarks=i,iangle=-90](p)(0,0){$p$}
  	\node[Nmarks=ir,iangle=-90](q)(20,0){$q$}
  	\node[Nmarks=r](r)(40,0){$r$}

	\drawloop(p){$a,b:0$}
  	\drawedge(p,q){$b:0$}
	\drawloop(q){$a:1$}
  	\drawedge(q,r){$b:0$}
	\drawloop(r){$a,b:0$}

  	\node[Nmarks=ir,iangle=180](s)(70,5){$s$}
	\drawloop[loopangle=90](s){$a:0$}
	\drawloop[loopangle=-90](s){$b:1$}
\end{picture}
\end{center}
\end{figure}

$$\tr{a} = 
\left(\begin{array}{cccc}
0 & \infty & \infty & \infty\\
\infty & \color{red}{1} & \infty & \infty\\
\infty & \infty & 0 & \infty\\
\infty & \infty & \infty & 0
\end{array}\right)
\qquad
\tr{a^n} = 
\left(\begin{array}{cccc}
0 & \infty & \infty & \infty\\
\infty & \color{red}{n} & \infty & \infty\\
\infty & \infty & 0 & \infty\\
\infty & \infty & \infty & 0
\end{array}\right)$$

$$\tr{a^\sharp} = 
\left(\begin{array}{cccc}
0 & \infty & \infty & \infty\\
\infty & \color{red}{\infty} & \infty & \infty\\
\infty & \infty & 0 & \infty\\
\infty & \infty & \infty & 0
\end{array}\right)
\qquad 
\only<2>{\left\{\begin{array}{cc}
k \in \NN & \textrm{best run has value } k\\
\infty & \textrm{arbitrarily high value}\\
\infty & \textrm{no run}
\end{array}\right.}$$
\end{frame}

\begin{frame}{Leung's algorithm}
To ensure termination we project the tropical semiring $(\NN \cup \infty, \min, +)$ 
into the \textbf{finite} semiring $(\set{0,1,\infty},\min,+)$.

\vskip2em
\pause

Compute a monoid inside the monoid $\MM_{Q \times Q}(\set{0,1,\infty}, \min, +)$.

\begin{itemize}
	\item Compute $\tr{a}$ for $a \in A$.
	\item Close under product and stabilization.
	\item If there exists a matrix $M$ such that $I \cdot M \cdot F = \infty$ then ``unlimited'',
	otherwise ``limited''.
\end{itemize}
\end{frame}

\begin{frame}{Leung's algorithm: termination and correction}

Termination: the monoid $\MM_{Q \times Q}(\set{0,1,\infty}, \min, +)$ is finite.

\vskip2em
\pause

Correction: the proof is complicated, and relies on Simon's theorem.
\end{frame}

\section{Towards an algebraic treatment of probabilistic automata}

\begin{frame}{Our objective (again)}
Decide the value $1$ problem for a \textit{subclass} of probabilistic automata,
by \textbf{algebraic} and \textbf{non-numerical} means.
\pause

\begin{itemize}
	\item \textbf{algebraic:} focus on the automaton structure,
	\item \textbf{non-numerical:} abstract away the values.
\end{itemize}

\pause
Hence we consider non-deterministic automata:
we project $(\RR,+,\cdot)$ into the boolean semiring $(\set{0,1},+,\cdot)$.
\pause
\begin{figure}
\begin{center}
\begin{picture}(20,15)(0,5)
	\gasset{Nw=6,Nh=6}

  	\node[Nmarks=i,iangle=180](1)(0,5){$1$}
  	\node[Nmarks=r](2)(20,5){$2$}

	\drawloop(1){$a$}
  	\drawedge(1,2){$a$}
	\drawloop(2){$a$}
\end{picture}
\end{center}
\end{figure}
\end{frame}

\subsection{First tries}

\begin{frame}{Defining stabilization}
\begin{figure}
\begin{center}
\begin{picture}(20,10)(0,0)
	\gasset{Nw=6,Nh=6}

  	\node[Nmarks=i,iangle=180](1)(0,5){$1$}
  	\node[Nmarks=r](2)(20,5){$2$}

	\drawloop(1){$a$}
  	\drawedge(1,2){$a$}
	\drawloop(2){$a$}
\end{picture}
\end{center}
\end{figure}

$$\tr{a} = 
\left(\begin{array}{ccc}
1 & 1 \\
0 & 1
\end{array}\right)
\qquad
\only<2,3>{
\tr{a^\sharp} = 
\left(\begin{array}{ccc}
0 & 1\\
0 & 1
\end{array}\right)}$$

In $\tr{a}$, the state $1$ is transient and the state $2$ is recurrent.

\only<3>{
$$M^\sharp(s,t) = 
\left\{\begin{array}{ll}
1 & \textrm{if } M(s,t) = 1 \textrm{ and } t \textrm{ recurrent in } M,\\
0 & \textrm{otherwise.}
\end{array}\right.$$

(This definition gives an asymmetric monoid, this is unusual.)}
\end{frame}

\begin{frame}{A first algorithm}
Compute a monoid inside the \textbf{finite} monoid $\MM_{Q \times Q}(\set{0,1},+,\cdot)$.

\begin{itemize}
	\item Compute $\tr{a}$ for $a \in A$:
$$\tr{a}(s,t) = 
\left\{\begin{array}{ll}
1 & \textrm{if } \prob{\AA}(s \xrightarrow{a} t) > 0,\\
0 & \textrm{otherwise.}
\end{array}\right.$$	
	\item Close under product and stabilization.
	\item \pause If there exists a matrix $M$ such that 
	$$\forall t \in Q, \quad M(s_0,t) = 1 \Rightarrow t \in F$$
	then ``$\AA$ has value $1$'',
	otherwise ``$\AA$ does not have value $1$''.
\end{itemize}
\end{frame}

\begin{frame}{An example}
\begin{figure}
\begin{center}
\begin{picture}(60,15)(0,-5)
	\gasset{Nw=6,Nh=6,loopdiam=6}

  	\node[Nmarks=i,iangle=-90](0)(20,5){$0$}
  	\node(1)(40,5){$1$}
  	\node[Nmarks=r](F)(60,5){$F$}
  	\node(bottom)(0,5){$\perp$}

	\drawloop(0){$a$}
	\drawloop(1){$a$}
	\drawloop[loopangle=0](F){$a,b$}
	\drawloop[loopangle=180](bottom){$a,b$}

  	\drawedge(0,bottom){$b$}
  	\drawedge[curvedepth=2](0,1){$a$}
  	\drawedge[curvedepth=2](1,0){$b$}
  	\drawedge(1,F){$b$}
\end{picture}
\pause
\begin{picture}(60,10)(0,0)
	\gasset{Nw=5,Nh=5,loopdiam=4}

	\put(-20,4){$\tr{a}$}

  	\node(0)(20,5){$0$}
  	\node(1)(40,5){$1$}
  	\node[Nmarks=r](F)(60,5){$F$}
  	\node(bottom)(0,5){$\perp$}

	\drawloop(0){}
	\drawloop(1){}
	\drawloop[loopangle=0](F){}
	\drawloop[loopangle=180](bottom){}

  	\drawedge(0,1){}
\end{picture}
\pause
\begin{picture}(60,10)(0,0)
	\gasset{Nw=5,Nh=5,loopdiam=4}

	\put(-20,4){$\tr{b}$}

  	\node(0)(20,5){$0$}
  	\node(1)(40,5){$1$}
  	\node[Nmarks=r](F)(60,5){$F$}
  	\node(bottom)(0,5){$\perp$}

	\drawloop[loopangle=0](F){}
	\drawloop[loopangle=180](bottom){}

  	\drawedge(0,bottom){}
  	\drawedge(1,0){}
  	\drawedge(1,F){}
\end{picture}
\pause
\begin{picture}(60,10)(0,0)
	\gasset{Nw=5,Nh=5,loopdiam=4}

	\put(-20,4){$\tr{a^\sharp}$}

  	\node(0)(20,5){$0$}
  	\node(1)(40,5){$1$}
  	\node[Nmarks=r](F)(60,5){$F$}
  	\node(bottom)(0,5){$\perp$}

	\drawloop(1){}
	\drawloop[loopangle=0](F){}
	\drawloop[loopangle=180](bottom){}

  	\drawedge(0,1){}
\end{picture}
\pause
\begin{picture}(60,10)(0,0)
	\gasset{Nw=5,Nh=5,loopdiam=4}

	\put(-22,4){$\tr{a^\sharp \cdot b}$}

  	\node(0)(20,5){$0$}
  	\node(1)(40,5){$1$}
  	\node[Nmarks=r](F)(60,5){$F$}
  	\node(bottom)(0,5){$\perp$}

	\drawloop(0){}
	\drawloop[loopangle=0](F){}
	\drawloop[loopangle=180](bottom){}

  	\drawedge[curvedepth=-4](0,F){}
  	\drawedge(1,0){}
  	\drawedge(1,F){}
\end{picture}
\pause
\begin{picture}(60,10)(0,0)
	\gasset{Nw=5,Nh=5,loopdiam=4}

	\put(-25,4){$\tr{(a^\sharp \cdot b)^\sharp}$}

  	\node(0)(20,5){$0$}
  	\node(1)(40,5){$1$}
  	\node[Nmarks=r](F)(60,5){$F$}
  	\node(bottom)(0,5){$\perp$}

	\drawloop[loopangle=0](F){}
	\drawloop[loopangle=180](bottom){}

  	\drawedge[curvedepth=-4](0,F){}
  	\drawedge(1,F){}
\end{picture}
\end{center}
\end{figure}
\end{frame}

\begin{frame}{Correctness}
\begin{theorem}
If there exists a matrix $M$ such that 
$$\forall t \in Q, \quad M(s_0,t) = 1 \Rightarrow t \in F$$
then $\AA$ has value $1$.
\end{theorem}
\pause
But the value $1$ problem is undecidable, so\ldots
\end{frame}

\begin{frame}{No completeness}
\begin{figure}
\begin{center}
\begin{picture}(60,35)(0,0)
	\gasset{Nw=7,Nh=7}

  	\node[Nmarks=i,iangle=-90](0)(30,15){$0$}
  	\node(L1)(0,10){$L_1$}
  	\node[Nmarks=r](L2)(0,30){$L_2$}
  	\node(R1)(60,10){$R_1$}
  	\node(R2)(60,30){$R_2$}

	\drawloop(0){$a$}

  	\drawedge[curvedepth=5,ELside=l](0,L1){$b,\frac{1}{2}$}
  	\drawedge[curvedepth=5,ELside=l](L1,0){$a,1-x$}
  	\drawedge(L1,L2){$b$}
	\drawloop[loopangle=-135](L1){$a,x$}
	\drawloop[loopangle=90](L2){$a,b$}

  	\drawedge[curvedepth=-5,ELside=r](0,R1){$b,\frac{1}{2}$}
  	\drawedge[curvedepth=-5,ELside=r](R1,0){$a,x$}
  	\drawedge[ELside=r](R1,R2){$b$}
	\drawloop[loopangle=-45](R1){$a,1-x$}
	\drawloop(R2){$a,b$}
\end{picture}
\end{center}
\end{figure}

%For all matrices $M$,
%$$\exists t \in Q, \quad M(s_0,t) = 1 \wedge t \notin F$$
%but $\AA$ has value $1$. 

Left and right parts are symmetric, so for all $M$:
$$M(0,L_2) = 1 \Longleftrightarrow M(0,R_2) = 1.$$
\end{frame}

\subsection{Leaks}

\begin{frame}{An example}
\begin{figure}
\begin{center}
\begin{picture}(50,30)(0,0)
	\gasset{Nw=6,Nh=6}

  	\node[Nmarks=i,iangle=0](L1)(15,15){$1$}
  	\node(L2)(0,30){$2$}
  	\node[Nmarks=r](L3)(0,0){$3$}

  	\drawedge(L1,L2){$b$}
  	\drawedge[curvedepth=-5,ELside=r](L1,L3){$a$}
  	\drawedge[curvedepth=-5,ELside=r](L3,L1){$b$}
	\drawloop(L1){$a$}
	\drawloop[loopangle=135](L2){$a,b$}
	\drawloop[loopangle=215](L3){$a$}

	\drawline[linewidth=.3,AHnb=0](26,-3)(26,40)

\pause

	\put(45,-10){$\tr{a^\sharp \cdot b}$}

  	\node(L1b)(55,15){$1$}
  	\node(L2b)(40,30){$2$}
  	\node[Nmarks=r](L3b)(40,0){$3$}

  	\drawedge(L3b,L1b){}
	\drawloop(L1b){}
	\drawloop[loopangle=135](L2b){}

\only<3>{\drawedge[dash={1}0](L1b,L2b){$\varepsilon$}}
\end{picture}
\end{center}
\end{figure}
\end{frame}

\subsection{The good semiring}

\begin{frame}{A three-valued semiring}
Instead of $(\set{0,1},+,\cdot)$ we compute in $(\set{0,\varepsilon,1},+,\cdot)$,
where $0 < \varepsilon < 1$.
\pause
$$\begin{array}{c||c|c|c}
+ & \quad 0 \quad & \quad \varepsilon \quad & \quad 1 \quad\\
\hline
0 & 0 & \varepsilon & 1 \\
\varepsilon & \varepsilon & \varepsilon & 1 \\
1 & 1 & 1 & 1 \\
\end{array}
\hspace{1cm}
\begin{array}{c||c|c|c}
\cdot & \quad 0 \quad & \quad \varepsilon \quad & \quad 1 \quad\\
\hline
0 & 0 & 0 & 0 \\
\varepsilon & 0 & \varepsilon & \varepsilon \\
1 & 0 & \varepsilon & 1 \\
\end{array}$$
\end{frame}

\begin{frame}{The algorithm}
\begin{itemize}
	\item Compute $\tr{a}$ for $a \in A$:
$$\tr{a}(s,t) = 
\left\{\begin{array}{ll}
1 & \textrm{if } \prob{\AA}(s \xrightarrow{a} t) > 0,\\
0 & \textrm{otherwise.}
\end{array}\right.$$	
	\item Close under product and stabilization:
$$M^\sharp(s,t) = 
\left\{\begin{array}{ll}
1 & \textrm{if } M(s,t) = 1 \textrm{ and } t \textrm{ recurrent in } M,\\
\varepsilon & \textrm{if } M(s,t) = 1 \textrm{ and } t \textrm{ transient in } M,\\
\varepsilon & \textrm{if } M(s,t) = \varepsilon,\\
0 & \textrm{otherwise.}
\end{array}\right.$$
	\item \pause If there exists a matrix $M$ such that 
	$$\forall t \in Q, \quad M(s_0,t) = 1 \Rightarrow t \in F$$
	then ``$\AA$ has value $1$'',
	otherwise ``$\AA$ does not have value $1$''.
\end{itemize}
\end{frame}

\begin{frame}{The control lemma}
We say that a word $w$ reify $M$ in $\MMAA$ if:
\begin{itemize}
	\item $M = \tr{a}$ and $w = a$;
	\item $M = M_1 \cdot M_2$ and there exists 
$w_1$ and $w_2$ reifying $M_1$ and $M_2$, respectively,
such that $w = w_1 \cdot w_2$;
	\item $M = N^\sharp$ and there exists $x_1,\ldots,x_n$ each reifying $N$,
such that $w = x_1 \ldots x_n$ for some $n \geq 1$.
\end{itemize}

\pause
\begin{lemma}[The control lemma]
For all $M$ in $\MMAA$, for all words $w$ reifying $M$, 
for all states $s,t$ in $Q$, we have:
$$M(s,t) \neq 0 \Longleftrightarrow 
\prob{\AA}(s \xrightarrow{w} t) > 0.$$
\end{lemma}
\end{frame}

\begin{frame}{Leaktight automata}
\begin{definition}
An automaton $\AA$ is leaktight if for all $M$,
we have
$$M(s,t) = \varepsilon \Longrightarrow (s \textrm{ is transient) or } (M(t,s) = 1).$$
\end{definition}
\pause

\begin{theorem}[Fijalkow,Gimbert,Oualhadj]
The value $1$ problem is decidable for leaktight automata.
\end{theorem}
\end{frame}

\subsection{The completeness proof using Simon's theorem}

\begin{frame}{Decomposition trees}

\begin{fact}
The set $\MMAA$ computed by the algorithm is a stabilization monoid.
\end{fact}

\begin{definition}
A \emph{decomposition tree} of a word $w \in A^+$ is a finite unranked ordered tree,
whose nodes have labels in $(A^+,\MMAA)$ and such that:
\begin{itemize}
	\item the root is labeled by $(w,u)$, for some $u \in \MMAA$,
	\item every leaf is labeled by $(a,\tr{a})$ where $a$ is a letter,
	\item every internal node with two children labeled by 
$(w_1,u_1)$ and $(w_2,u_2)$ is labeled by 
$(w_1 \cdot w_2,u_1 \cdot u_2)$,
	\item for every internal node with three or more children, 
there exists $e \in E(M)$ such that 
the node is labeled by $(w_1 \ldots w_n,e^\sharp)$
and its children are labeled by $(w_1,e),\ldots,(w_n,e)$.
\end{itemize}
\end{definition}
\end{frame}

\begin{frame}{Bounding the height of a decomposition tree}
In a decomposition tree,
an iteration node is said discontinuous if $M^\sharp \neq M$.
The span of a decomposition tree is the maximal length of a path that contains
no discontinuous path.

\pause
\begin{theorem}[Simon, 1990]
Every word $w \in A^+$ has a decomposition tree whose span is less than 
$3 \cdot |\MMAA|$.
\end{theorem}

\pause
\begin{lemma}[Simon, 1990]
Let $M \in E(\MMAA)$, if $M^\sharp \neq M$, then $M^\sharp <_\JJ M$.
\end{lemma}

\pause
\begin{corollary}
Every word $w \in A^+$ has a decomposition tree whose height is less than 
$3 \cdot |\MMAA| \cdot J(\AA)$.
\end{corollary}
\end{frame}

\begin{frame}{Bounding the acceptance probability from below}

\begin{lemma}
There exists a positive rational number $\eta$ which depends only on $\AA$ such that:
for all words $w \in A^+$, there exists $M$ in $\MMAA$
satisfying

for all states $s,t \in Q$,
$$M(s,t) = 1 \Rightarrow \prob{\AA}(s \xrightarrow{w} t) \geq \eta.$$
\end{lemma}

\pause
Proof idea: given $w$, consider a decomposition tree of bounded height,
and prove by induction that the lower bound $2^{-h+1}$ holds at depth $h$,
going from leaves to the root.
\end{frame}

\begin{frame}{The case of an iteration node (1)}

The node is labelled by $(w_1 \ldots w_n,\tr{u^\sharp})$
and its children are labelled by $(w_1,\tr{u}),\ldots,(w_n,\tr{u})$,
where $\tr{u}$ is idempotent, and $\eta$ a lower bound shared by the $n \geq 3$ children.
%We have for all $i$, $s,t$ in $Q$, if $\tr{u}(s,t) = 1$ then 
%$\prob{\AA}(s \xrightarrow{w_i} t) \geq \eta$: call this property $(\dagger)$.

\pause
Let $s,t$ such that $\tr{u^\sharp}(s,t) = 1$, then:
$$\prob{\AA}(s \xrightarrow{w_1 \ldots w_n} t)
\geq \underbrace{\prob{\AA}(s \xrightarrow{w_1} t)}_{\geq \eta} 
\cdot \underbrace{\prob{\AA}(t \xrightarrow{w_2 \ldots w_n} t)}_{\geq \eta}
\geq \eta^2.$$

\pause
The left inequality follows from induction hypothesis,
since $\tr{u}(s,t) = 1$.
\end{frame}

\begin{frame}{The case of an iteration node (2)}
Consider the right inequality: $\prob{\AA}(t \xrightarrow{w_2 \ldots w_n} t) \geq \eta$

Let $C = \set{q \mid \tr{u}(t,q) \neq 0}$, we have:
$$\begin{array}{lll}
\prob{\AA}(t \xrightarrow{w_2 \ldots w_n} t) 
& = & \sum_{q \in C} \prob{\AA}(t \xrightarrow{w_2 \ldots w_{n-1}} q) 
\cdot \underbrace{\prob{\AA}(q \xrightarrow{w_n} t)}_{\geq \eta} \\
& \geq & \eta \cdot 
\underbrace{\sum_{q \in C} \prob{\AA}(t \xrightarrow{w_2 \ldots w_{n-1}} q)}_{= 1} = \eta
\end{array}$$
Indeed, since $t$ is recurrent and thanks to the leaktight assumption,
we have $C \subseteq \set{q \mid \tr{u}(q,t) = 1}$,
so the inequality follows from induction hypothesis,
and the equality from the ``control lemma''.
\end{frame}

\begin{frame}{What I didn't (and won't) say}
\begin{itemize}
	\item One can decide whether an automaton is leaktight in $\PSPACE$,
	\item The value $1$ problem for probabilistic leaktight automata is $\PSPACE$-complete,
	\item The class of leaktight automata subsumes all subclasses of probabilistic automata
whose value $1$ problem is known to be decidable,
	\item The class of leaktight automata is closed under parallel composition and synchronized product.
\end{itemize}
\end{frame}

\begin{frame}{The end.}
Thanks for your attention!

\begin{center}
\includegraphics[scale=0.2]{mario.eps}
\end{center}
\end{frame}

\end{document}